\scanpage{37}

phismes de restriction
\[
H^n(\overline{X},\overline{F})\qquad H^n(\overline{Y},\overline{\unsure{H|B}}\ (\overline{F}))\qquad\text{pour }n\quad0\ ,
\]
et les homomorphismes de changement de base
\[
h\ (R^n\ (\overline{F}))\qquad R^n\quad{}_{Y}\ (\overline{h}\ (\overline{F}))\qquad\text{pour }n\quad0,
\]
soient bijectifs. Alors pour tout faisceau abélien de torsion $F$ sur $X$, les homomorphismes de restriction
\[
H^n(X,F)\qquad H^n(Y,h\ (F))\qquad\text{pour }n\quad0
\]
sont \deleted{des} bijectifs.
% The spaces between the formula members, and between n and 0, contain no visible signs.

\uline{Démonstration}.\quad (i)\quad Pour l'assertion\deleted{n} non respée, on utilise le carré commutatif (où on pose $\overline{F}=\quad(F)$):
% No arrow strokes are visible in this square.
\diagram{
\begin{tikzcd}[ampersand replacement=\&,column sep=3em,row sep=3em]
H^0(X,F) \& H^0(Y,F) \\
H^0(\overline{X},\overline{F}) \& H^0(\overline{Y},\overline{F})\quad,
\end{tikzcd}
}
où par hypothèse la deuxième flèche horizontale est injective, ce qui montre que deux éléments de $H^0(X\deleted{n},F)$ ayant même image dans $H^0(Y,F)$ ont aussi même image dans $H^0(\overline{X},\overline{F})$, \deleted{donc}\inserted{i.e.} ont mêmes images dans les fibres \deleted{géo} de $F$ en les points géométriques au dessus des points de $(\overline{X})$, donc dans \uline{toutes} les fibres puisque $X=Y\quad(\overline{X})$, ce qui prouve qu'elles sont égales. Pour prouver l'assertion respée, soit $F$ un faisceau d'ensembles sur $X$ et prouvons que $H^0(X,F)\qquad H^0(Y,h\ (F))$ est surjectif (donc bijecttif d'après ce qui précède). Montrons qu'il suffit pour ceci de trouver une injection $F\qquad G$, telle que l'assertion analogue soit vraie pour $G$. En effet, d'après les propriétés d'exactitude générales d'une catégorie de faisceaux, $F$ peut alors s'écrire comme un conoyau d'un couple de flèches $G\qquad H$ (où on peut prendre pour $H$ la somme amalgamée $\unsure{\deleted{F}G|\deleted{G}F}\quad{}_{F}G$, si on veut). Tenant compte que le foncteur $H$ est exact\deleted{n}, \deleted{on obtient via} et
% archive sheet number: 36
